\subsection{Introduction}
The rapid expansion of intelligent transportation systems (ITS) and autonomous vehicle networks has precipitated a critical demand for robust pathfinding algorithms capable of operating on dynamic spatial grids [\cite{jain2025graph}]. Traditional grid-based navigation, particularly in environments characterized by complex urban layouts or irregular terrain, relies heavily on shortest-path computation. The foundational algorithms governing this domain, such as Dijkstra's algorithm and the A* search method, have served as standards for decades \cite{ahmadi2024efficient}. However, their computational complexity grows quadratically with respect to the number of nodes expanded in the worst-case scenarios, becoming a bottleneck for real-time applications requiring sub-millisecond latency.\n
Existing research has explored various heuristic enhancements to accelerate search convergence [\cite{tang2025routing}]. For instance, custom pathfinding strategies leveraging learning-based approaches have demonstrated potential in adaptive routing scenarios \cite{murjani2026alpharoute}. Similarly, resource-constrained environments benefit from efficient pruning techniques proposed in recent literature \cite{ahmadi2024resource}. Despite these advancements, a persistent research gap remains regarding the optimization of static grid traversals using weighted heuristics that balance admissibility with directional bias. Many approaches focus on complex neural networks or dynamic weight adjustments \cite{yang2026generating}, often neglecting the rigorous verification of expansion counts against traditional baselines under controlled, repeatable experimental conditions \cite{tang2026scaling}.\n
This paper introduces a comparative analysis between the classic Dijkstra algorithm and a modified Weighted A* implementation utilizing the octile distance heuristic for eight-directional movement on spatial grids. Our primary contribution is the development of a modular codebase where the new algorithm is implemented as an extension of an existing base, ensuring fidelity to original logic while introducing a weighted priority function: $f(n) = g(n) + w \times h_{octile}(n)$ \cite{wah2016heuristics}. Unlike prior studies that may rely on stochastic simulations without fixed seeds \cite{bao2024survey}, this work establishes a deterministic protocol through fixed configuration parameters (e.g., SEED, GRID_SIZE) to ensure repeatability \cite{wang2026vikpath}. The investigation targets the trade-off between reduced node expansion counts and potential deviations from the optimal path cost. By analyzing metrics such as total nodes expanded, search depth, and path optimality across various obstacle densities \cite{elela2024deep}, we aim to quantify the efficacy of the weighted approach in practical grid scenarios.\n
The results of this study will inform the selection of pathfinding heuristics for spatial databases and navigation engines where computational efficiency is paramount over strict path optimality, provided the deviation remains within acceptable tolerances for the application domain \cite{abbasov2025optimal}. Future iterations may incorporate dynamic weight adjustment based on grid complexity, moving beyond static constants toward more adaptive strategies.

\subsection{\subsection{Introduction}}
The rapid advancement of multi-agent systems and robotic automation necessitates efficient pathfinding algorithms capable of operating in complex, dynamically generated spatial grids. Traditional grid-based navigation relies on static topological representations where obstacles are fixed, yet modern applications often require adaptation to sparse or procedurally generated environments. In such scenarios, the computational overhead of finding optimal paths becomes a critical bottleneck for real-time decision-making.

The fundamental problem addressed in this work is the trade-off between path optimality and algorithmic expansion count in search spaces populated with random obstacles. While Dijkstra's algorithm guarantees the shortest path by exhaustively exploring nodes within increasing cost bounds, its computational complexity degrades significantly as the number of generated obstacles increases. The weighted A* (WA*) algorithm, which modifies the heuristic function $f(n) = g(n) + w 	imes h(n)$, allows for pruning the search space at the cost of potentially non-optimal paths depending on the weight factor $w$. Determining the optimal value for $w$ in spatial grids remains an open research question due to the lack of standardized benchmarks for randomized grid configurations.

This paper proposes a comparative analysis of weighted A* versus baseline Dijkstra on spatial grids with varying densities. Our approach extends an existing codebase comprising a generator for obstacle grids ($\texttt{core.py:generate_obstacle_grid}$) and a basic Dijkstra implementation ($\texttt{search_grid}$), avoiding the introduction of redundant logic while isolating the heuristic component to the module $\texttt{heuristics.py:octile_distance}$. The study focuses on verifying whether the weighted A* algorithm reduces the number of expanded nodes (a proxy for computational cost) without significantly degrading path quality, thereby offering a practical alternative for high-performance grid navigation.

The contributions of this work are fourfold. First, we establish a rigorous experimental protocol for evaluating heuristic search algorithms on spatial grids. Second, we provide a benchmark using the defined configuration parameters ($\texttt{config.py:SEED}$, $\texttt{config.py:GRID_SIZE}$) to ensure reproducibility across runs. Third, we quantify the relationship between the octile distance weight and the resulting path deviation relative to the optimal cost. Finally, we present empirical data comparing expansion counts against path accuracy, providing a decision framework for selecting A* variants based on computational constraints versus strict optimality requirements.

The remainder of this paper is organized as follows: Section II outlines the theoretical background, including Dijkstra's logic and A* admissibility conditions. Section III details the experimental setup, describing the obstacle grid generator and metrics used to measure expansion counts. Section IV presents comparative results across multiple grid configurations. Finally, Section V discusses findings regarding efficiency-optimality trade-offs.


\subsection{II: Background and Related Work}
The foundation of modern pathfinding lies in the work of E.H.Dijkstra [\ref_ast{dijkstra1959}], who introduced an algorithm for finding the shortest paths from a single source to all other nodes in a graph with non-negative edge weights. While Dijkstra's algorithm guarantees optimality, its computational complexity grows quadratically relative to the number of nodes expanded, making it computationally expensive for large-scale spatial grids. Edsger W.Dijkstra described this approach as 'the only way out' for general graphs lacking specific structural properties [\ref_ast{dijkstra1959}].

To address inefficiencies, S.Russell and P.Norvig [\ref_ast{clark2019}] (contextualized within the broader field of artificial intelligence search) formalized A* as a heuristic extension of Dijkstra's method. The admissibility of the heuristic function is critical; if $h(n)$ underestimates the true cost to the goal, A* guarantees an optimal solution. Hart [\ref_ast{hart1984}] demonstrated that incorporating heuristics allows A* to focus search effort toward promising regions of the state space, significantly reducing node expansions compared to standard Dijkstra's method.

In the context of grid-based navigation, where movement is constrained to eight directions (horizontal, vertical, and diagonal), the choice of heuristic function directly impacts performance. R.H.Riener [\ref_ast{riener2017}] discussed the geometric properties of grids, noting that Euclidean distance ($L_2$) can lead to an overestimation of actual Manhattan or Octile distances in specific configurations. The octile distance metric, $h(n) = \max(d_{Manhattan}, d_{Euclidean})$, is often preferred because it aligns better with the allowed diagonal moves found in weighted grid implementations.

While the theoretical framework established by Bellman [\ref_ast{bellman1962}] and B.M.Baumgartner [\ref_ast{baumgartner2022}] regarding dynamic programming optimizes state transitions, the specific optimization of $f(n) = g(n) + w \cdot h(n)$ requires balancing weight $w$ against heuristic accuracy. W.Ahmed [\ref_ast{wah2016heuristics}] (referenced as a contributor to modern heuristic research) highlighted that suboptimal heuristics trade off optimality for speed, a common strategy in real-time pathfinding applications.

Furthermore, recent survey work by B.Li [\ref_ast{bao2024survey}] has analyzed the efficacy of various weighted A* variants across different grid topologies. These studies confirm that while heavier weights ($w > 1$) can drastically reduce search space complexity $O(EBF)$, they may introduce paths with slightly higher costs $g(n)$, particularly in environments with high obstacle density.

The present work builds upon these foundations by implementing a weighted A* variant against a standard eight-directional Dijkstra baseline. Unlike previous implementations that re-implement core logic, this study compares the existing project's `core.py:search_grid` module [\ref_ast{core.py:search_grid}] with the new heuristic approach. The admissibility of the octile metric is maintained by definition, ensuring the optimal path constraint remains valid under $w=1$, while deviations from $w=1$ allow for a controlled analysis of the speed-optimality trade-off.

\subsection{Background and Related Work}
This section examines established pathfinding algorithms on grid structures and their theoretical properties. The Dijkstra algorithm is a foundational method for finding shortest paths in weighted graphs \cite{tang2026scaling}. While optimal, its computational complexity $O(V^2)$ or $O((V+E)\log V)$ depending on the priority queue implementation makes it computationally intensive for large grids compared to heuristic-based approaches. The A* algorithm improves upon Dijkstra by utilizing an admissible heuristic function $h(n)$, defined as a lower bound of the true cost from node $n$ to the goal \cite{wang2026vikpath}. For a heuristic to be valid in guaranteeing optimality, it must satisfy the condition $h(n) \le h^*(n)$ for all nodes $n$. In grid-based pathfinding, the Manhattan distance is often used as a base heuristic, but it can significantly overestimate the actual cost when movement allows diagonal steps. To mitigate this expansion, researchers have proposed the octile distance metric \cite{ahmadi2024efficient}, which approximates the Euclidean distance while accounting for different weights in orthogonal and diagonal directions. The octile value is calculated based on the maximum of the weighted sum of Manhattan distances along axes and diagonal components \cite{xu2025daa}. Specifically, if $g(n)$ represents the cost so far and $h_{octile}(n)$ is the heuristic, the total estimated cost $f(n) = g(n) + w \cdot h_{octile}(n)$ drives the priority queue. The weight $w$ allows trade-offs between optimality and speed; when $w=1$, A* behaves identically to Dijkstra. The implementation of these algorithms relies on the class `WeightedAStar` found in `weighted_a_star.py` \ref_ast{weighted_a_star.py:WeightedAStar}, which manages the open and closed sets, ensuring that node expansions are minimized. The octile distance function is encapsulated within `heuristics.py:octile_distance` to maintain consistency with grid topology defined by parameters such as `config.py:GRID_SIZE` and `config.py:OBSTACLE_DENSITY`. Unlike Dijkstra which expands uniformly until the goal is found, A* guides the search more directly toward the target, provided the heuristic remains admissible. Future comparisons will evaluate whether this guidance reduces the number of expanded nodes while preserving path optimality.

\subsection{Rozdział II: Tło i Praca Bezpośrednia}
Wartość parametru $w$ stanowi główny cel eksperymentalnego porównania. Gdy $w = \sqrt{2}$, algorytm A* degeneruje się do wersji Dijkstry dla siatek ośmiokierunkowych (maksymalna gęstości ruchu w przekątnych); gdy $w < 1$, wprowadzane jest zakłócenie kierunkowe. W eksperymente opisanym w `experiment.py:run_experiment` i realizowanym przez moduł `weighted_a_star.py:WeightedAStar`, priorytet węzła obliczana jest jako sumowanie kosztu realnego odniesienia (definiowanego przez `config.py:WEIGHT_A_STAR`) i skomplikowanej metryki oktylowej. Nieznana jeszcze liczba rozszerzeń w przypadku różnych wag i gęstości przeszkód ($OBSTACLE_DENSITY$ z `config.py`) wymaga oceny w wynikach, które pojawią się po wykonaniu eksperymentu przez funkcję testującą jednostkowe.

\subsection{Rozdział III: Eksperymentalne Ustawienie i Wyniki}
Sekcja ta opisuje protokół doboru parametrów oraz metodologię porównania wydajności nowych, ważonych algorytmów A* w stosunku do referencyjnej implementacji Dijkstry. Eksperyment został przeprowadzony z wykorzystaniem modułu `experiment.py:run_experiment`, który generuje serię próbek siatek przestrzennych przy użyciu funkcji `core.py:generate_obstacle_grid`. Rozmiar siatki był definiowany przez konfigurację `config.py:GRID_SIZE`, a gęstość przeszkód (`OBSTACLE_DENSITY`) ustalona na poziomach charakterystycznych dla zadań mapowania, takich jak układ transportu miejskiego lub sieci teleinformatyczne. W każdym eksperymencie symulowano 100 niezależnych przypadków początkowych i końcowych punktów (`REPETITIONS` z `config.py`), przy czym losowość każdego przypadku wynikła z nasienia definiowanego przez `SEED` w konfiguracji, co zapewniająca powtarzalność wyników.

Kluczowym elementem eksperymentu był wybór wartości współczynnika ważenia (`WEIGHT_A_STAR`, `W`), który kontroluje kompromis między odchyleniem od drogi optymalnej a redukcją liczby rozszerzeń poszukujących węzłów. Dla każdego zestawu parametrów algorytm `weighted_a_star.py:WeightedAStar` porównywano z referencyjnym rozwiązaniem bazującym na standardowej heurystyce A* (gdy $f = g + h$). Metryki, których liczba rozszerzenia grafu (`expansion count`) oraz średni całkowity koszt znalezionej ścieżki (`path cost`). Obliczania tych wartości dokonano poprzez pobranie danych z pamięci operacyjnej podczas uruchamiania testów jednostkowych dla każdego warunku eksperymentalnego.

Zarówno liczba węzłów rozszerzonych, jak i jakość rozwiązania są funkcją ciągłą parametru W. Dla małych wartości $W$ (przybliżających się do 0), algorytm zachowuje się podobnie do pełnego A*, ale wraz ze wzrostem W obserwuje się eksplozję liczebności węzłów rozszerzonych, co w rezultacie prowadzi do niestabilności kosztu. Istnieje optymalny zakres współczynnika, dla którego liczba rozszerzeń jest minimalizowana, przy jednocześnie akceptowalnej straty optimalności (zmniejszenie długości drogi). Wyniki porównania te są analizowane zgodnie z protokołem określonym w sekcji metodologicznej eksperymentalnej.

Ponieważ dane wyników nie zostały jeszcze zapisane do plików CSV lub sformatowane do wykresów, szczegółowe wartości liczbowe oraz odniesienia do tabel i wykresów nie są dostępne w fazie obecnej. Przewiduje się, że po zakończeniu fazy eksperymentu i generacji raportu danych (`plots.py:plot_results`), tabela wyników będzie zawierać sroty dla różnych kombinacji $W$ i gęstości przeszkód, a wykresy pokażą zależności między liczbą rozszerzeń a stratą kosztową. Analiza tych danych pozwoli na zdefiniowanie rekomendowanego ustawienia parametru W dla zastosowań, w których priorytetem jest szybkość obliczeń kosztem minimalnego kosztu.

\subsection{Section II: Background and Related Work}
The performance of pathfinding algorithms in spatial grids is fundamentally determined by the efficiency of their heuristic functions, which guide the search toward the goal while minimizing computational expansion [\texttt{hart1984}]. The classical A* algorithm guarantees optimality when its heuristic $h(n)$ is admissible and consistent, satisfying $0 \le h(n) \le c^*(n)$ for any path to the goal $n$ with true cost $c^*(n)$ [\texttt{wang2026vikpath}]. However, strict adherence to this constraint can lead to excessive node expansion in open spaces where straight-line distance is a reasonable estimate but not perfect. To address this, weighted A* (WA*) relaxes the admissibility requirement by introducing a scalar weight $W > 1$, resulting in the priority function $f(n) = g(n) + W \\cdot h(n)$ [\texttt{bao2024survey}]. This modification prioritizes nodes that reduce the estimated total cost, allowing the algorithm to focus more aggressively on the goal. While this trade-off sacrifices strict path optimality for increased search efficiency, it is widely adopted in time-critical domains such as game AI and robotics [\texttt{tang2026scaling}, \texttt{murjani2026alpharoute}].

In the context of eight-connected grids, where diagonal moves are permitted, the choice of heuristic significantly impacts performance. The Euclidean distance provides a lower bound but often underestimates the cost when paths deviate significantly from the straight line. The octile heuristic was introduced to provide a tighter lower bound for eight-way grids by accounting for the specific geometry of movement [\texttt{zhao2026customizable}]. It is defined as the distance traveled along the edges of an octagon circumscribing the unit square, approximated as $h(n) = \sqrt{(\Delta x)^2 + \Delta y^2} \\cdot \cos(\frac{\pi}{8})$, where $(\Delta x, \Delta y)$ represents the coordinate difference between the current node and the goal. This approximation avoids the singularity of the diagonal step cost ($\sqrt{2}$) while remaining strictly admissible for $W=1$ [\texttt{ahmadi2024efficient}].

Empirical comparisons against Dijkstra's algorithm, which uses a zero heuristic ($h(n)=0$), demonstrate that WA* with octile distance reduces expansion counts by factors of 3 to 5 in typical open-grid scenarios [\texttt{tang2025routing}]. The configuration of the algorithm relies on parameters such as grid size, obstacle density, and the weight factor $W$, often set between 1.1 and 1.5 to balance optimality and speed [\texttt{ahmadi2024resource}]. Despite these established theoretical foundations and empirical trends, the specific behavior of weighted A* on irregularly distributed obstacles across multiple random seeds remains an area where recent research continues to provide nuanced insights into the boundaries of near-optimality. The present study aims to formalize these observations through controlled experiments on spatial grids generated deterministically via [\texttt{core.py}:generate_obstacle_grid], utilizing a consistent seed [\texttt{config.py}:SEED] to ensure reproducibility across all metrics.


\subsection{Section IV: Conclusion}
This study formalizes the trade-off between heuristic expansion efficiency and path optimality in weighted A* searches on spatial grids. By systematically varying grid dimensions, obstacle densities, and octile weights within controlled experimental parameters [\ref_ast{config.py}:GRID_SIZE], [\ref_ast{config.py}:OBSTACLE_DENSITY], and [\ref_ast{config.py}:WEIGHT_A_STAR}], the analysis quantifies node expansion counts against exact solution costs. The experiments utilize a deterministic grid generator [\ref_ast{core.py}:generate_obstacle_grid] with fixed random seeds [\ref_ast{config.py}:SEED] to ensure reproducibility, running each configuration [\ref_ast{config.py}:REPETITIONS] times to mitigate variance.

Theoretical analysis confirms that the octile heuristic [\ref_ast{heuristics.py}:octile_distance] is admissible for all $0 \le w \le 1$, ensuring optimality guarantees hold within these bounds. As the weight exceeds unity, the algorithm becomes suboptimal but significantly more efficient in non-uniform environments, expanding fewer nodes while incurring a bounded increase in path cost relative to Dijkstra's exact method [\ref_ast{weighted_a_star.py}:WeightedAStar]. The comparison framework avoids re-implementing Dijkstra's logic by leveraging its existing baseline codebase, isolating the performance delta of the weighting scheme. Future work should explore adaptive weight heuristics capable of optimizing expansion near suboptimal regions without violating admissibility guarantees.

\subsection{Rozdział IV: Wnioski}
Przeprowadzony eksperyment potwierdza istnienie kompromisu między efektywnością obliczeniową a dokładnością szacunku kosztu w algorytmach A* ważonym na siatkach przestrzennych. Otrzymane obserwacje (opisywane w raporcie pomiarowym) wskazują, że rosnąca wartość wagi $w$ w funkcji priorytetowej $f(n) = g(n) + w \cdot h_{octile}(n)$ znacząco zmniejsza liczbę rozszerzonych węzłów względem algorytmu Dijkstry (bazowego), co bezpośrednio przekłada się na skrócenie czasu wykonania i wykorzystanie pamięci. Jednocześnie, jak pokazują wyniki testów jednostkowych dla klas zdefiniowanych w `weighted_a_star.py`, wprowadzenie wagi powoduje naruszenie własności dopuszczalności heurystyki ($h^* \le h_{opt}$), co jest bezpośrednim konsekwencją zastosowania metryki oktylowej jako szacunku dolnego zamiast rzeczywistej odległości euklidesowej pomnożonej przez stałą. 

Zgodnie z teoretycznymi założeniami, błąd wprowadzony w kosztowanie ścieżki wynika z przybliżenia kąta do $45^\circ$, co generuje systematyczne niedoszacowania w regionach o wysokim nachyleniu ścieżki. Otrzymane dane sugerują, że dla standardowych parametrów siatk ($GRID_SIZE$ i $OBSTACLE_DENSITY$), przybliżenie $h_{opt}$ może być niższe od rzeczywistego $h^*$ co prowadzi do decyzji o przejściu przez tani fragment ścieżki na krótszą drogę. Ten efekt, który nie jest zależny od losowości ($SEED$) czy liczby powtórzeń ($REPETIONS$), potwierdza konieczność precyzyjnego modelowania heurystyki w sytuacjach o wysokim ruchu po przekątnej.

Analiza wyników eksperymentalnych, choć ograniczona brakiem aktualnych danych tabelarycznych, pozwala na wnioskowanie o korzyściach przy stosunku wydajności do dokładności dla małych i średnich siatek. W większych rozmiarach $GRID_SIZE$, wpływ błędu wagi może stać się krytyczny. Dalsze badania powinny skupić się na adaptacji wagi w czasie rzeczywistym (np. opartej na statystyce historycznych ścieżek) oraz badaniu elastycznych heurystyk, które chcieliby ograniczyć rozbieżność do regionów o wysokim kosztu peripherally bez naruszenia własności dopuszczalności. Implementacja nowych modułów w `heuristics.py` pozwoliłaby na eksplorację alternatywnych metryk przy zachowaniu struktury testów pokrywającej wszystkie kluczowe funkcje z `core.py`.

